\clearpage\selectlanguage{english}\fancyhead[L]{\small GENERAL INVESTIGATIONS}
\SourceLeaf{134}
\begin{center}
GENERAL INVESTIGATIONS

\medskip
OF THE INFINITE SERIES
\[
\begin{aligned}&1+\frac{\alpha\beta}{1\mathbin{.}\gamma}x\\
&{}+\frac{\alpha(\alpha+1)\beta(\beta+1)}{1\mathbin{.}2\mathbin{.}\gamma(\gamma+1)}xx\\
&{}+\frac{\alpha(\alpha+1)(\alpha+2)\beta(\beta+1)(\beta+2)}{1\mathbin{.}2\mathbin{.}3\mathbin{.}\gamma(\gamma+1)(\gamma+2)}x^3+\mathrm{etc.}
\end{aligned}
\]
FIRST PART.

\medskip
\rule{0.13\linewidth}{0.4pt}

\medskip
\textit{INTRODUCTION.}

\end{center}
\SixthArticle{1}
The series that we undertake to investigate in this paper can be regarded as a function of four quantities $\alpha,\beta,\gamma,x$, which we shall call its \textit{elements}, distinguishing them in order as the first element $\alpha$, the second $\beta$, the third $\gamma$ and the fourth $x$. Clearly, the first element can be interchanged with the second. Thus, if for brevity we denote our series by $F(\alpha,\beta,\gamma,x)$, we shall have $F(\beta,\alpha,\gamma,x)=F(\alpha,\beta,\gamma,x)$.

\SixthArticle{2}
When definite values are assigned to the elements $\alpha,\beta,\gamma$, our series becomes a function of the single variable $x$. It clearly terminates after the $1-\alpha^{\mathrm{th}}$ or $1-\beta^{\mathrm{th}}$ term if $\alpha-1$ or $\beta-1$ is a negative integer, but in the remaining cases it continues to infinity. In the former case the series represents a rational algebraic function; in the latter it usually represents a transcendental function. The third element $\gamma$ must be neither a negative integer nor $=0$, so that we do not encounter infinitely large terms.
\SourceLeaf{135}
\SixthArticle{3}
The coefficients of the powers $x^m$, $x^{m+1}$ in our series are in the ratio
\[
1+\frac{\gamma+1}{m}+\frac{\gamma}{mm}:1+\frac{\alpha+\beta}{m}+\frac{\alpha\beta}{mm}
\]
and therefore their ratio comes closer to unity as $m$ is taken larger. Thus, if a definite value is also assigned to the fourth element $x$, convergence or divergence will depend on its nature. Whenever $x$ is assigned a real value, positive or negative, smaller than unity, the series will certainly converge, if not from the very beginning, then after a certain interval, and will lead to a finite, completely determined sum. The same will happen for an imaginary value of $x$ of the form $a+b\BareRoot{-1}$ whenever $aa+bb<1$. Conversely, for a real value of $x$ greater than unity, or an imaginary value of the form $a+b\BareRoot{-1}$ whenever $aa+bb>1$, the series will necessarily diverge, if not at once, then after a certain interval, so that there can be no question of its \textit{sum}. Finally, for the value $x=1$ (or, more generally, a value of the form $a+b\BareRoot{-1}$ whenever $aa+bb=1$), the convergence or divergence of the series will depend on the nature of $\alpha,\beta,\gamma$ themselves. We shall discuss this, and in particular the sum of the series for $x=1$, in the third Section.

It is therefore clear that, insofar as our function is defined as the sum of a series, the investigation is by its nature restricted to those cases in which the series actually converges. Thus it is a senseless question to ask what the value of the series is for a value of $x$ greater than unity. Below, however, from the fourth Section onward, we shall base our function on a higher principle that allows the most general application.

\SixthArticle{4}
Differentiating our series, with only the fourth element $x$ regarded as variable, leads to a similar function, since clearly
\[
\frac{\mathrm{d}F(\alpha,\beta,\gamma,x)}{\mathrm{d}x}
=\frac{\alpha\beta}{\gamma}F(\alpha+1,\beta+1,\gamma+1,x)
\]
The same holds for repeated differentiation.

\SixthArticle{5}
It will be worthwhile to set down here some functions that can be reduced to our series and are used very frequently throughout analysis.
\SourceLeaf{136}
\[
\begin{array}{rl}
\mathrm{I.}&(t+u)^n=t^nF(-n,\beta,\beta,-\frac{u}{t})
\end{array}
\]
where the element $\beta$ is arbitrary.
\[
\begin{array}{rl}
\mathrm{II.}&(t+u)^n+(t-u)^n=2t^nF(-\frac{1}{2}n,-\frac{1}{2}n+\frac{1}{2},\frac{1}{2},\frac{uu}{tt})
\end{array}
\]
\[
\begin{array}{rl}
\mathrm{III.}&(t+u)^n+t^n=2t^nF(-n,\omega,2\omega,-\frac{u}{t})
\end{array}
\]
where $\omega$ denotes an infinitely small quantity.
\[
\begin{array}{rl}
\mathrm{IV.}&(t+u)^n-(t-u)^n=2nt^{n-1}uF(-\frac{1}{2}n+\frac{1}{2},-\frac{1}{2}n+1,\frac{3}{2},\frac{uu}{tt})
\end{array}
\]
\[
\begin{array}{rl}
\mathrm{V.}&(t+u)^n-t^n=nt^{n-1}uF(1-n,1.\ 2,-\frac{u}{t})
\end{array}
\]
\[
\begin{array}{rl}
\mathrm{VI.}&\log(1+t)=tF(1,1,2,-t)
\end{array}
\]
\[
\begin{array}{rl}
\mathrm{VII.}&\log\frac{1+t}{1-t}=2tF(\frac{1}{2},1,\frac{3}{2},tt)
\end{array}
\]
\Needspace{6\baselineskip}
\[
\begin{array}{rl}
\mathrm{VIII.}&e^t=F(1,k,1,\frac{t}{k})=1+tF(1,k,2,\frac{t}{k})=1+t+\frac{1}{2}ttF(1,k,3,\frac{t}{k})\ \mbox{etc.}
\end{array}
\]
where $e$ denotes the base of the natural logarithms and $k$ an infinitely large number.
\[
\begin{array}{rl}
\mathrm{IX.}&e^t+e^{-t}=2F(k,k',\frac{1}{2},\frac{tt}{4kk'})
\end{array}
\]
where $k$, $k'$ denote infinitely large numbers.
\[
\begin{array}{rl}
\mathrm{X.}&e^t-e^{-t}=2tF(k,k',\frac{3}{2},\frac{tt}{4kk'})
\end{array}
\]
\[
\begin{array}{rl}
\mathrm{XI.}&\sin t=tF(k,k',\frac{3}{2},-\frac{tt}{4kk'})
\end{array}
\]
\[
\begin{array}{rl}
\mathrm{XII.}&\cos t=F(k,k',\frac{1}{2},-\frac{tt}{4kk'})
\end{array}
\]
\[
\begin{array}{rl}
\mathrm{XIII.}&t=\sin t.\ F(\frac{1}{2},\frac{1}{2},\frac{3}{2},\sin t^2)
\end{array}
\]
\[
\begin{array}{rl}
\mathrm{XIV.}&t=\sin t.\ \cos t.\ F(1,1,\frac{3}{2},\sin t^2)
\end{array}
\]
\[
\begin{array}{rl}
\mathrm{XV.}&t=\mathrm{tang}\,t.\ F(\frac{1}{2},1,\frac{3}{2},-\mathrm{tang}\,t^2)
\end{array}
\]
\[
\begin{array}{rl}
\mathrm{XVI.}&\sin nt=n\sin t.\ F(\frac{1}{2}n+\frac{1}{2},-\frac{1}{2}n+\frac{1}{2},\frac{3}{2},\sin t^2)
\end{array}
\]
\[
\begin{array}{rl}
\mathrm{XVII.}&\sin nt=n\sin t.\ \cos t.\ F(\frac{1}{2}n+1,-\frac{1}{2}n+1,\frac{3}{2},\sin t^2)
\end{array}
\]
\[
\begin{array}{rl}
\mathrm{XVIII.}&\sin nt=n\sin t.\ \cos t^{n-1}F(-\frac{1}{2}n+1,-\frac{1}{2}n+\frac{1}{2},\frac{3}{2},-\mathrm{tang}\,t^2)
\end{array}
\]
\[
\begin{array}{rl}
\mathrm{XIX.}&\sin nt=n\sin t.\ \cos t^{-n-1}F(\frac{1}{2}n+1,\frac{1}{2}n+\frac{1}{2},\frac{3}{2},-\mathrm{tang}\,t^2)
\end{array}
\]
\[
\begin{array}{rl}
\mathrm{XX.}&\cos nt=F(\frac{1}{2}n,-\frac{1}{2}n,\frac{1}{2},\sin t^2)
\end{array}
\]
\[
\begin{array}{rl}
\mathrm{XXI.}&\cos nt=\cos t.\ F(\frac{1}{2}n+\frac{1}{2},-\frac{1}{2}n+\frac{1}{2},\frac{1}{2},\sin t^2)
\end{array}
\]
\[
\begin{array}{rl}
\mathrm{XXII.}&\cos nt=\cos t^nF(-\frac{1}{2}n,-\frac{1}{2}n+\frac{1}{2},\frac{1}{2},-\mathrm{tang}\,t^2)
\end{array}
\]
\[
\begin{array}{rl}
\mathrm{XXIII.}&\cos nt=\cos t^{-n}F(\frac{1}{2}n+\frac{1}{2},\frac{1}{2}n,\frac{1}{2},-\mathrm{tang}\,t^2)
\end{array}
\]
